Test-time training cho một mô hình ngôn ngữ khả năng tự cải thiện ngay tại
thời điểm suy luận, trên một bài toán khó duy nhất. Nhưng self-distillation
theo kiểu on-policy rất dễ giậm chân tại chỗ: nếu mô hình không thể tự sinh
ra một lời giải đúng, thì cũng chẳng có gì đúng để nó học theo. Khóa luận
này tìm cách tăng tốc test-time discovery trong sinh mã phức tạp bằng
\emph{execution-guided self-distillation}: dùng execution feedback
(test case thất bại, lỗi runtime) làm privileged context dẫn dắt từng
bước cập nhật. Tôi đề xuất một cách tổ chức \emph{teacher-first}:
thay vì distill chính rollout thất bại của student, self-teacher sinh
trước các quỹ đạo dưới feedback; một verifier và một judge lọc chúng theo
tính đúng và tính độc lập; rồi student mới được distill về phía những quỹ
đạo đã qua lọc đó. Cách tổ chức này đặt phương pháp đề xuất ở vị trí giao
thoa giữa SDPO on-policy và SFT off-policy. Trong một đánh giá đối chiếu
(matched evaluation) quy mô trung bình trên các bài toán LiveCodeBench khó
bằng mô hình 4B, teacher-first vượt trội đáng kể so với baseline
student-first (kiểm định Wilcoxon signed-rank). Đặc biệt trên những bài khó
nhất, phương pháp này thoát được flat-reward trap, cạm bẫy vốn luôn
khiến student-first giậm chân tại mức 0. Một so sánh cân bằng tính toán
(compute-matched) cho thấy lợi thế này không phải là hệ quả ảo sinh ra từ
khối lượng lấy mẫu: ở cùng một ngân sách generation, teacher-first vẫn
chiến thắng và đạt được xác suất khám phá (discovery) mục tiêu với số
lượng generation ít hơn đáng kể. Khi kích hoạt cơ chế suy nghĩ (thinking),
tôi tiến hành đo lường epistemic verbalization tại thời điểm suy luận và
phát hiện ra rằng: quá trình chưng cất (distill) từ một feedback-conditioned
context làm suy giảm các uncertainty marker, một minh chứng trực tiếp
cho hiện tượng suy giảm mà Kim và cộng sự đã mô tả. Cuối cùng, việc đánh
giá pipeline trên tác vụ suy luận toán học (AIME 2026) đã vạch rõ một ranh
giới áp dụng rất nghiêm ngặt (rigid domain boundary): khi feedback chỉ gói
gọn trong một giá trị số tĩnh thay vì một quy trình thuật toán, teacher
policy sẽ suy biến thành hành vi sao chép đáp án nông cạn mà không mang
lại bất kỳ sự cải thiện năng lực nào, qua đó thiết lập nên một ràng buộc
khắt khe giữa quy trình và giá trị (procedure-versus-value) đối với
self-distillation tại thời điểm suy luận.

\vspace{0.5cm}
\textbf{Từ khóa:} test-time training, self-distillation, sinh mã,
execution feedback, discovery@k, epistemic suppression, ranh giới miền.
